\documentclass[11pt]{article}
\usepackage[margin=25mm]{geometry}
\usepackage{amsmath,amssymb,amsthm}
\usepackage[hidelinks]{hyperref}
\newtheorem{theorem}{Theorem}
\newcommand{\R}{\mathbb R}
\newcommand{\Pp}{\mathcal P}
\title{Conjecture 146: the stated normalization has no weak probability limit}
\author{Local development pilot; not an upstream submission}
\date{8 October 2026}
\begin{document}
\maketitle
\noindent The source uses primes $p\le N$, denominator $N$, and scaling $N$.
We preserve all three. Its supremum domain and law of $\alpha$ are unstated.
We use the conventional Kolmogorov--Smirnov domain $[0,1]$ and prove a stronger
statement covering every probability law of $\alpha$, even one changing with $N$.
A supremum over all real thresholds would be unbounded; that interpretation is
not used as a shortcut.

For $N\ge1$, put $\Pp_N=\{p\in\mathbb N:p\text{ prime},\ p\le N\}$ and
\[
 F_N(\alpha,t)=\frac{\#\{p\in\Pp_N:\{\alpha p\}<t\}}{N},\qquad
 D_N(\alpha)=\sup_{0\le t\le1}|F_N(\alpha,t)-t|,\qquad
 X_N(\alpha)=N D_N(\alpha).
\]
The Lean definitions are total at $N=0$ using real division by zero; every
limiting argument only uses $N\ge4$.

\begin{theorem}
For any sequence of Borel probability measures $\mu_N$ on $\R$, the actual
pushforward laws $\lambda_N=(X_N)_*\mu_N$ have no weak limit which is a
probability measure on $\R$. In particular they cannot converge to the law of
any real-valued Brownian-bridge functional.
\end{theorem}

\begin{proof}
Every fractional part belongs to $[0,1)$, so
$F_N(\alpha,1)=\#\Pp_N/N$ for every $\alpha$. Hence
\[
 X_N(\alpha)\ge N-\#\Pp_N\ge N/2-2\quad(N\ge2).
\tag{1}
\]
For the last inequality, all primes other than $2$ are odd. Counting those
integers gives $\#\Pp_N\le N/2+2$; no prime number theorem is needed. The Lean
proof uses the corresponding arithmetic prime-counting bound.

These are genuine measurable statistics. For fixed $t$, each comparison
$\{\alpha p\}<t$ is measurable: multiplication and fractional part are Borel
measurable. A finite sum of indicators therefore makes $F_N(\cdot,t)$
measurable. Furthermore,
\[
 D_N(\alpha)=\sup_{q\in\mathbb Q\cap[0,1]}
                   |F_N(\alpha,q)-q|.
\tag{2}
\]
To check the nontrivial direction, take $t>0$. The finite collection of samples
strictly below $t$ has a maximum below $t$, if nonempty. Choose rational $q<t$
arbitrarily close to $t$ and above every such sample. Then the strict-threshold
count at $q$ equals the count at $t$, and the displayed errors approach one
another. The case $t=0$ is itself rational. This argument also handles $t=1$,
repeated samples and an empty collection. The ranges of errors are nonempty
and bounded, as checked in Lean. Equation (2) is a countable supremum of
measurable functions; thus $D_N$ and $X_N$ are measurable, and $\lambda_N$ is
indeed a probability measure.

Use the bounded continuous, everywhere strictly positive test
\[
 f(x)=\frac1{1+|x|},\qquad 0<f(x)\le1.
\]
For $N\ge4$, (1) implies
\[
 0\le\int f\,d\lambda_N
   =\int f(X_N(\alpha))\,d\mu_N(\alpha)
   \le \frac1{N/2-1}\longrightarrow0.
\tag{3}
\]
The equality is the proved pushforward integral identity. If $\lambda_N$
converged weakly to a real probability measure $\nu$, this bounded continuous
test would also give $\int f\,d\lambda_N\to\int f\,d\nu$.
Yet $\int f\,d\nu>0$: a nonnegative measurable function has zero integral only
if it vanishes almost everywhere, whereas $f$ is positive everywhere and
$\nu(\R)=1$. This contradicts (3).
\end{proof}

\paragraph{Formal correspondence.}
\texttt{Route2.primeKS} is the actual supremum over the subtype $[0,1]$;
\texttt{scaledKS} multiplies it by $N$.
\texttt{measurable\_primeKS} proves (2)'s measurability consequence, and
\texttt{laws\_probability} and \texttt{test\_transport} certify the actual image
laws. The zero-limit and strictly-positive-integral theorems supply the
contradictory limits.
The final theorem is \texttt{Route2.no\_probability\_weak\_limit}.
Its explicit weak-convergence predicate tests all continuous real functions
with values in $[0,1]$, using their nonnegative lower integrals. This is the
usual normalized bounded-continuous-test formulation for probabilities; in
particular standard weak convergence necessarily passes the test used above.
No extra premise assumes measurability, the escape bound, or an object bridge.

\paragraph{Provenance and limits.}
The endpoint bound, rational-supremum equality and fixed-threshold CDF
measurability reuse disclosed Phase~1c B partial work; an uncompiled Phase~1c A
draft supplied the proposed integral route. This round completes the
measurability, probability-law and weak-limit connections. Those prior ideas
are not claimed as new discoveries. The exact original, pinned Lean sources,
source correspondence and independent verification records accompany this
local package. Local verification does not establish submission eligibility,
a first solve, or upstream acceptance.
\end{document}
